\chapter{Komunikacja Sieciowa: Retrofit (Wykład 10)}

W poprzednim rozdziale nauczyliśmy się przechowywać dane lokalnie na urządzeniu. Jest to podejście typu \textbf{Offline-First}. Jednak współczesne aplikacje mobilne rzadko działają w izolacji. Aby pobrać prognozę pogody, wyświetlić listę najnowszych wiadomości czy zsynchronizować zadania między urządzeniami, musimy skomunikować się z zewnętrznym serwerem.

Komunikacja sieciowa w aplikacjach mobilnych opiera się zazwyczaj na architekturze REST (Representational State Transfer).
\begin{itemize}
    \item \textbf{Klient (Twoja Aplikacja):} Wysyła żądanie HTTP (Request) pod określony adres URL.
    \item \textbf{Serwer (API):} Przetwarza żądanie i odsyła odpowiedź (Response).
\end{itemize}

Operacja sieciowa jest zjawiskiem wysoce nieprzewidywalnym. Czas odpowiedzi serwera (Latency) zależy od jakości zasięgu (LTE/5G/WiFi), obciążenia łącza czy odległości geograficznej od serwera. Może to trwać od 50 ms do nawet kilku sekund.

Ze względu na tę nieprzewidywalność, system Android narzuca rygorystyczne ograniczenie: \textbf{Wykonywanie jakichkolwiek operacji sieciowych na Wątku Głównym (UI Thread) jest zabronione.}

Próba wykonania kodu sieciowego wewnątrz texttt{MainActivity} (bez użycia korutyn lub wątków tła) zakończy się natychmiastowym wyrzuceniem wyjątku \texttt{NetworkOnMainThreadException}.

\section{Uprawnienia aplikacji}

Domyślnie aplikacja na Androida nie ma prawa łączyć się z Internetem. Aby to umożliwić, musimy zadeklarować odpowiednie uprawnienie w pliku manifestu. Jest to uprawnienie typu \textit{Normal}, co oznacza, że jest przyznawane automatycznie przy instalacji i nie wymaga odrębnej zgody użytkownika w czasie działania (Runtime Permission).

\begin{minted}[frame=single,framesep=10pt]{xml}
<manifest xmlns:android="http://schemas.android.com/apk/res/android">

    <uses-permission android:name="android.permission.INTERNET" />

    <application ...>
        ...
    </application>
</manifest>
\end{minted}

\section{Biblioteka Retrofit 2}

W początkach systemu Android do obsługi sieci używano klasy \texttt{HttpURLConnection}. Wymagało to ręcznego otwierania strumieni, buforowania danych, obsługi kodowania znaków i ręcznego parsowania tekstu zwracanego formatu (np. JSON) na obiekty Javy.

Firma Square stworzyła bibliotekę \texttt{Retrofit}, która zrewolucjonizowała ten proces. Retrofit to tzw. \textit{Type-safe HTTP client}, który pozwala programiście zdefiniować API jako prosty interfejs w Kotlinie.
Zamiast pisać kod implementujący połączenie, my tylko \textit{deklarujemy}, jak wygląda zapytanie (metoda HTTP, ścieżka, parametry). Retrofit automatycznie generuje kod wykonawczy.

Główne zalety:
\begin{enumerate}
    \item \textbf{Abstrakcja:} Ukrywa niskopoziomowe szczegóły protokołu HTTP.
    \item \textbf{Automatyczna konwersja (Serializacja/Deserializacja):} Współpracuje z bibliotekami takimi jak GSON, Moshi czy Scalars, automatycznie zamieniając odebraną odpowiedź na obiekty \texttt{data class}.
    \item \textbf{Integracja z Kotlin Coroutines:} Retrofit natywnie wspiera funkcje zawieszające (\texttt{suspend}).
\end{enumerate}

Dzięki wsparciu dla korutyn, zapytanie sieciowe w kodzie wygląda jak zwykłe wywołanie synchroniczne, ale zachowuje pełne bezpieczeństwo wątkowe (Main-Safety).

\begin{minted}[frame=single,framesep=10pt]{kotlin}
// Definicja interfejsu (Retrofit)
interface NewsApi {
    @GET("articles")
    suspend fun getArticles(): List<Article>
}

// Użycie w ViewModelu
viewModelScope.launch {
    try {
        // To wywołanie:
        // 1. ZAWIESI korutynę (zwolni wątek UI).
        // 2. Wykona zapytanie sieciowe w tle.
        // 3. WZNOWI korutynę na wątku UI, gdy nadejdą dane.
        val articles = api.getArticles()
        
        // Tutaj mamy już gotową listę obiektów, możemy zaktualizować UI
        uiState = articles
    } catch (e: Exception) {
        // Obsługa błędów (np. brak internetu)
    }
}
\end{minted}

W powyższym przykładzie nie musimy używać \texttt{withContext(Dispatchers.IO)}. Przełączanie odbywa się, podobnie jak w Room i DataStore, automatycznie.

\section{Metody protokołu HTTP (Semantyka)}

W architekturze REST każda operacja na danych (tzw. CRUD) jest ściśle powiązana z odpowiednią metodą protokołu HTTP. Wybór metody informuje serwer o intencjach klienta.

Najczęściej używane metody to:

\begin{itemize}
    \item \textbf{GET (Pobierz):} Służy wyłącznie do pobierania danych. Jest operacją bezpieczną i zazwyczaj nie powinna modyfikować stanu serwera (np. pobierz listę artykułów, pobierz szczegóły użytkownika).
    \item \textbf{POST (Wyślij/Utwórz):} Służy do wysyłania nowych danych na serwer w celu utworzenia zasobu (np. dodaj nowy komentarz, zarejestruj użytkownika). Dane przesyłane są w ciele żądania (Body).
    \item \textbf{PUT (Zastąp):} Służy do pełnej aktualizacji istniejącego zasobu. Klient wysyła kompletną nową wersję obiektu, która zastępuje starą.
    \item \textbf{PATCH (Łataj):} Służy do częściowej modyfikacji zasobu. Wysyłamy tylko te pola, które mają ulec zmianie (np. zmiana samego hasła, bez przesyłania reszty profilu).
    \item \textbf{DELETE (Usuń):} Jak nazwa wskazuje, usuwa wskazany zasób z serwera.
\end{itemize}

\begin{table}[h]
\centering
\begin{tabular}{|l|l|l|}
\hline
\textbf{Operacja CRUD} & \textbf{Metoda HTTP} & \textbf{Adnotacja Retrofit} \\ \hline
Create (Utwórz)        & POST                 & \texttt{@POST}              \\ \hline
Read (Odczytaj)        & GET                  & \texttt{@GET}               \\ \hline
Update (Aktualizuj)    & PUT / PATCH          & \texttt{@PUT} / \texttt{@PATCH} \\ \hline
Delete (Usuń)          & DELETE               & \texttt{@DELETE}            \\ \hline
\end{tabular}
\caption{Mapowanie operacji CRUD na metody HTTP.}
\end{table}

\section{Modelowanie Danych (Data Transfer Objects)}

Zanim zaczniemy pobierać dane, musimy przygotować dla nich \textit{pojemniki}. W inżynierii oprogramowania obiekty służące wyłącznie do przesyłania danych między podsystemami (np. Serwer $\rightarrow$ Aplikacja) nazywamy \textbf{DTO (Data Transfer Objects)}.

Serwery REST API najczęściej zwracają dane w formacie JSON (JavaScript Object Notation). Naszym zadaniem jest stworzenie klas w Kotlinie (\texttt{data class}), które odwzorowują strukturę tego JSON-a (Dodajmy tutaj że ten proces można zautomatyzować przez instalację odpowiednich pluginów).

Standardy nazewnictwa w JSON i Kotlinie często się różnią:
\begin{itemize}
    \item \textbf{JSON:} Często używa notacji \texttt{snake\_case} (np. \texttt{first\_name}, \texttt{created\_at}).
    \item \textbf{Kotlin:} Wymaga notacji \texttt{camelCase} (np. \texttt{firstName}, \texttt{createdAt}).
\end{itemize}

Aby rozwiązać ten konflikt bez łamania konwencji języka, używamy adnotacji \texttt{@SerializedName} (z biblioteki GSON). Działa ona jak \textit{mapa}, mówiąca konwerterowi: \textit{"Pole 'user\_avatar\_url' z JSON-a przypisz do zmiennej 'avatarUrl'"}.

Załóżmy, że serwer zwraca następujący obiekt artykułu:

\begin{minted}[frame=single,framesep=10pt]{json}
{
  "id": 90210,
  "title": "Nowy Android 15 wydany!",
  "published_at": "2024-05-12T10:00:00Z",
  "author": {
    "name": "Jan Kowalski",
    "is_admin": true
  }
}
\end{minted}

Odpowiadające mu klasy w Kotlinie wyglądają następująco:

\begin{minted}[frame=single,framesep=10pt]{kotlin}
data class ArticleDto(
    // Nazwa identyczna jak w JSON -> adnotacja niepotrzebna
    val id: Long,
    
    val title: String,
    
    // Mapowanie snake_case -> camelCase
    @SerializedName("published_at")
    val publishedAt: String,
    
    // Obiekty zagnieżdżone
    val author: AuthorDto
)

data class AuthorDto(
    val name: String,
    
    @SerializedName("is_admin")
    val isAdmin: Boolean
)
\end{minted}

\section{Definiowanie Interfejsu API (Service)}

Sercem biblioteki Retrofit jest interfejs. To tutaj deklarujemy, jakie operacje nasza aplikacja może wykonać na serwerze. Nie piszemy ciała metod - robimy to wyłącznie za pomocą adnotacji i typów zwracanych. Retrofit wygeneruje kod wykonawczy w czasie działania aplikacji.

Najprostszy przypadek to pobranie zasobu ze stałego adresu URL.

\begin{minted}[frame=single,framesep=10pt]{kotlin}
interface NewsApiService {
    // Żądanie zostanie wysłane na adres: BASE_URL + "articles"
    // Funkcja MUSI być suspend, aby współpracować z korutynami.
    @GET("articles")
    suspend fun getArticles(): List<ArticleDto>
}
\end{minted}

Rzadko zdarza się, że pobieramy zawsze to samo. Często chcemy pobrać konkretny artykuł lub przefiltrować listę.

Parametr Path (@Path):
Używamy, gdy zmienna jest częścią ścieżki URL.
\begin{minted}[frame=single,framesep=10pt,fontsize=\footnotesize]{kotlin}
// Zapytanie o konkretny artykuł, np. GET articles/123
@GET("articles/{id}")
suspend fun getArticleDetail(
    @Path("id") articleId: Long
): ArticleDto
\end{minted}

Parametr Query (@Query):
Używamy do filtrowania, sortowania lub stronicowania (to, co występuje po znaku \texttt{?}).
\begin{minted}[frame=single,framesep=10pt]{kotlin}
// Zapytanie: GET articles?category=tech&sort=desc
@GET("articles")
suspend fun getArticles(
    @Query("category") category: String,
    @Query("sort") sortBy: String = "desc"
): List<ArticleDto>
\end{minted}

Gdy chcemy wysłać obiekt do serwera (np. utworzyć nowy komentarz), używamy metody POST oraz adnotacji \texttt{@Body}. Retrofit automatycznie zserializuje obiekt do formatu JSON przed wysłaniem.

\begin{minted}[frame=single,framesep=10pt]{kotlin}
data class CommentRequest(
    val articleId: Long,
    val content: String
)

interface NewsApiService {
    @POST("comments")
    suspend fun postComment(@Body request: CommentRequest): Response<Unit>
}
\end{minted}

Często API wymaga autoryzacji (np. klucza API). Możemy dodać nagłówek statycznie lub dynamicznie.

\begin{minted}[frame=single,framesep=10pt]{kotlin}
// 1. Statyczny nagłówek dla metody
@Headers("User-Agent: MyAndroidApp/1.0")
@GET("articles")
suspend fun getNews(): List<ArticleDto>

// 2. Dynamiczny nagłówek (np. token sesji)
@GET("profile")
suspend fun getProfile(
    @Header("Authorization") token: String
): UserDto
\end{minted}

\section{Konfiguracja Klienta (Retrofit Builder)}

Mając zdefiniowany model danych (DTO) oraz interfejs API, musimy utworzyć instancję biblioteki Retrofit. Proces ten przypomina konfigurację bazy danych Room - również tutaj zastosujemy wzorzec Singleton, aby nie tworzyć kosztownych połączeń sieciowych wielokrotnie.

\subsection{Zależności (Gradle)}

Najpierw upewnijmy się, że w pliku \texttt{build.gradle (Module: app)} znajdują się wymagane biblioteki. Retrofit jest modułowy - sam w sobie obsługuje tylko połączenia HTTP. Do parsowania JSON-a potrzebujemy dodatkowego \textit{konwertera}.

\begin{minted}[frame=single,framesep=10pt]{kotlin}
dependencies {
    // Silnik Retrofit
    implementation("com.squareup.retrofit2:retrofit:2.9.0")
    // Konwerter JSON -> Kotlin Object (GSON)
    implementation("com.squareup.retrofit2:converter-gson:2.9.0")
    
    // Opcjonalnie: OkHttp do logowania zapytań w konsoli (Debugging)
    implementation("com.squareup.okhttp3:logging-interceptor:4.11.0")
}
\end{minted}

\subsection{Implementacja Singletona}

Tworzymy obiekt, który będzie zarządzał konfiguracją. Kluczowe elementy to:
\begin{itemize}
    \item \textbf{Base URL:} Adres główny API. \textbf{Uwaga:} Musi zawsze kończyć się znakiem ukośnika \texttt{/}.
    \item \textbf{ConverterFactory:} Klasa odpowiedzialna za serializację/deserializację (tutaj GSON).
    \item \textbf{OkHttpClient:} Klient HTTP, w którym możemy skonfigurować np. timeouty (czas oczekiwania na serwer).
\end{itemize}

\begin{minted}[frame=single,framesep=10pt]{kotlin}
object RetrofitInstance {
    
    private const val BASE_URL = "https://newsapi.org/v2/"

    // Konfiguracja klienta HTTP (opcjonalne logowanie)
    private val client = OkHttpClient.Builder()
        .addInterceptor { chain ->
            val request = chain.request().newBuilder()
                // Możemy tu dodać globalny nagłówek, np. klucz API
                .addHeader("X-Api-Key", "TWOJ_KLUCZ_API") 
                .build()
            chain.proceed(request)
        }
        .build()

    // Leniwa inicjalizacja Retrofita
    val api: NewsApiService by lazy {
        Retrofit.Builder()
            .baseUrl(BASE_URL)
            .client(client)
            .addConverterFactory(GsonConverterFactory.create())
            .build()
            .create(NewsApiService::class.java)
    }
}
\end{minted}

\section{Obsługa Błędów i Stany UI}

W przeciwieństwie do lokalnej bazy danych, zapytanie sieciowe może zakończyć się niepowodzeniem na wiele sposobów (brak zasięgu, błąd 404, błąd 500). Aplikacja musi być na to gotowa.

Zamiast trzymać w ViewModelu osobne zmienne \texttt{isLoading}, \texttt{errorText}, \texttt{data}, eleganckim rozwiązaniem jest użycie \textbf{Sealed Interface}. Dzięki temu widok zawsze znajduje się w jednym, ściśle określonym stanie.

\begin{minted}[frame=single,framesep=10pt]{kotlin}
sealed interface NewsUiState {
    // 1. Stan początkowy / Ładowanie
    data object Loading : NewsUiState
    
    // 2. Sukces - mamy dane do wyświetlenia
    data class Success(val articles: List<ArticleDto>) : NewsUiState
    
    // 3. Błąd - coś poszło nie tak
    data class Error(val message: String) : NewsUiState
}
\end{minted}

\begin{minted}[frame=single,framesep=10pt]{kotlin}
class NewsViewModel : ViewModel() {

    // Stan wewnętrzny (edytowalny)
    private val _uiState = MutableStateFlow<NewsUiState>(NewsUiState.Loading)
    
    // Stan publiczny (tylko do odczytu)
    val uiState = _uiState.asStateFlow()

    init {
        fetchNews()
    }

    fun fetchNews() {
        viewModelScope.launch {
            // Ustawiamy stan na ładowanie
            _uiState.update { NewsUiState.Loading }
            
            try {
                // Wywołanie sieciowe (Retrofit + Suspend)
                val response = RetrofitInstance.api.getTopHeadlines()
                
                // Sukces: Aktualizujemy stan danymi
                _uiState.update { 
                    NewsUiState.Success(response.articles) 
                }
            } catch (e: IOException) {
                // Błąd sieci (brak internetu, timeout)
                _uiState.update { 
                    NewsUiState.Error("Brak połączenia z internetem") 
                }
            } catch (e: HttpException) {
                // Błąd serwera (4xx, 5xx)
                _uiState.update { 
                    NewsUiState.Error("Błąd serwera: ${e.code()}") 
                }
            } catch (e: Exception) {
                // Inne błędy
                _uiState.update { 
                    NewsUiState.Error("Nieznany błąd: ${e.message}") 
                }
            }
        }
    }
}
\end{minted}

\section{Przykład Praktyczny: Aplikacja News Reader}

Zepnijmy wszystko w całość. Stworzymy ekran, który reaguje na zmiany stanu: wyświetla kręciołek podczas ładowania, listę po pobraniu danych lub komunikat błędu w razie awarii.

\subsection{Model i API}

Zdefiniujmy strukturę danych i endpoint (jak w poprzednich punktach).

\begin{minted}[frame=single,framesep=10pt]{kotlin}
// DTO (Data Transfer Object)
data class NewsResponse(val articles: List<ArticleDto>)
data class ArticleDto(val title: String, val description: String?)

// Interfejs Retrofit
interface NewsApiService {
    @GET("top-headlines?country=pl")
    suspend fun getTopHeadlines(): NewsResponse
}
\end{minted}

\subsection{Warstwa Repozytorium}

Repozytorium ukrywa źródło danych. Jeśli w przyszłości dodamy cache w bazie Room, zmienimy kod tylko tutaj, a ViewModel pozostanie bez zmian.

\begin{minted}[frame=single,framesep=10pt]{kotlin}
class NewsRepository(private val apiService: NewsApiService) {
    
    // Funkcja zawieszająca, przekazująca żądanie do Retrofita
    suspend fun getNews(): List<ArticleDto> {
        return apiService.getTopHeadlines().articles
    }
}
\end{minted}

\subsection{ViewModel z obsługą StateFlow}

ViewModel przyjmuje repozytorium w konstruktorze. Używamy \texttt{StateFlow} do zarządzania stanem.

\begin{minted}[frame=single,framesep=10pt]{kotlin}
class NewsViewModel(private val repository: NewsRepository) : ViewModel() {

    // Backing property
    private val _uiState = MutableStateFlow<NewsUiState>(NewsUiState.Loading)
    val uiState = _uiState.asStateFlow()

    init {
        fetchNews()
    }

    fun fetchNews() {
        viewModelScope.launch {
            _uiState.update { NewsUiState.Loading }
            
            try {
                // Pobieramy dane z REPOZYTORIUM, nie bezpośrednio z API
                val articles = repository.getNews()
                
                _uiState.update { NewsUiState.Success(articles) }
            } catch (e: IOException) {
                _uiState.update { NewsUiState.Error("Brak internetu") }
            } catch (e: HttpException) {
                _uiState.update { NewsUiState.Error("Błąd serwera: ${e.code()}") }
            } catch (e: Exception) {
                _uiState.update { NewsUiState.Error("Nieznany błąd") }
            }
        }
    }
}

// Fabryka potrzebna do wstrzyknięcia Repozytorium do ViewModelu
class NewsViewModelFactory(private val repository: NewsRepository) : ViewModelProvider.Factory {
    override fun <T : ViewModel> create(modelClass: Class<T>): T {
        if (modelClass.isAssignableFrom(NewsViewModel::class.java)) {
            @Suppress("UNCHECKED_CAST")
            return NewsViewModel(repository) as T
        }
        throw IllegalArgumentException("Unknown ViewModel class")
    }
}
\end{minted}

\subsection{Widok (Compose)}

Warstwa UI obserwuje stan i reaguje na niego.

\begin{minted}[frame=single,framesep=10pt]{kotlin}
@Composable
fun NewsScreen(viewModel: NewsViewModel) { // ViewModel przekazywany jako parametr
    
    val state by viewModel.uiState.collectAsStateWithLifecycle()

    Scaffold(
        topBar = { TopAppBar(title = { Text("Wiadomości") }) }
    ) { padding ->
        Box(modifier = Modifier.padding(padding).fillMaxSize(), 
          contentAlignment = Alignment.Center) {
            when (val currentState = state) {
                is NewsUiState.Loading -> CircularProgressIndicator()
                is NewsUiState.Error -> {
                    Column(horizontalAlignment = Alignment.CenterHorizontally) {
                        Text(text = currentState.message)
                        Button(onClick = { viewModel.fetchNews() }) { Text("Ponów") }
                    }
                }
                is NewsUiState.Success -> {
                    LazyColumn(modifier = Modifier.fillMaxSize()) {
                        items(currentState.articles) { article ->
                            ArticleItem(article)
                        }
                    }
                }
            }
        }
    }
}
\end{minted}

\subsection{Punkt wejścia (MainActivity)}

\begin{minted}[frame=single,framesep=10pt,fontsize=\footnotesize]{kotlin}
class MainActivity : ComponentActivity() {
    override fun onCreate(savedInstanceState: Bundle?) {
        super.onCreate(savedInstanceState)
        
        // 1. Utworzenie instancji API (z Singletona)
        val apiService = RetrofitInstance.api
        
        // 2. Utworzenie Repozytorium
        val repository = NewsRepository(apiService)
        
        // 3. Utworzenie ViewModelu przez Fabrykę
        val viewModelFactory = NewsViewModelFactory(repository)

        setContent {
            val viewModel: NewsViewModel = viewModel(factory = viewModelFactory)
            NewsScreen(viewModel = viewModel)
        }
    }
}
\end{minted}

\section{Singleton vs Application Class}

Uważny czytelnik może zauważyć różnicę w sposobie inicjalizacji klienta sieciowego względem bazy danych z Rozdziału 9.
\begin{itemize}
    \item W przypadku \textbf{Room}, musieliśmy użyć klasy \texttt{Application}, ponieważ baza danych wymaga dostępu do \texttt{Context} (systemu plików urządzenia).
    \item W przypadku \textbf{Retrofita}, zastosowaliśmy uproszczenie w postaci statycznego obiektu (\texttt{object RetrofitInstance}). Jest to możliwe, ponieważ klient HTTP w swojej podstawowej formie nie zależy od systemu Android i nie wymaga \texttt{Contextu}.
\end{itemize}